%%%%%%%%%%%%%%%%%%%%%%%%%%%%%%%%%%%%%%%%%
% Important note:
% This template requires the resume.cls file to be in the same directory as the
% .tex file. The resume.cls file provides the resume style used for structuring the
% document.
%
%%%%%%%%%%%%%%%%%%%%%%%%%%%%%%%%%%%%%%%%%

%----------------------------------------------------------------------------------------
%	PACKAGES AND OTHER DOCUMENT CONFIGURATIONS
%----------------------------------------------------------------------------------------

\documentclass{resume} % Use the custom resume.cls style

\geometry{left=0.75in,top=0.6in,right=0.75in,bottom=0.6in} % Document margins

\newcommand{\tab}[1]{\hspace{.2667\textwidth}\rlap{#1}}
\newcommand{\itab}[1]{\hspace{0em}\rlap{#1}}
\name{Qijia Fan} % Your name
\address{\faEnvelope\ \href{mailto:fqj25@mails.tsinghua.edu.cn}{fqj25@mails.tsinghua.edu.cn} ~~|~~ \faGithub \ \href{https://github.com/qweryy0566}{qweryy0566}} % Your phone number and email

\begin{document}

%----------------------------------------------------------------------------------------
%	EDUCATION SECTION
%----------------------------------------------------------------------------------------

\begin{rSection}{Education and Research Experience}
%--copy and paste this region  if you need more--
{\bf Tsinghua University} \hfill {\em Sept. 2025 -- Present}
\\ \emph{Ph.D. student of Computer Science, advised by Prof. Yifan Song.}

{\bf Brown University} \hfill {\em Aug. 2024 -- Dec. 2024}
\\ \emph{Research Intern.}
% \begin{itemize}[nosep, leftmargin=*]
%   \item Proposed a fuzzy labeled private set intersection protocol with lower communication cost.
%   \item Analyzed the Hamming-distance distribution of face embeddings and selected suitable parameters for the hash table used in the protocol.
% \end{itemize}

{\bf Shanghai Jiao Tong University} \hfill {\em Aug. 2021 -- Jun. 2025} 
\\ \emph{Bachelor in Computer Science.}
\vspace{-0.6em}
\begin{itemize}[nosep, leftmargin=*]
  \item A member of {\bf ACM Honors Class}, an elite CS program in SJTU.
  \item GPA: 94.27/100, 4.11/4.3 ~|~ Ranking: 2/34
\end{itemize}

%--copy and paste this region  if you need more--

\end{rSection}
%----------------------------------------------------------------------------------------
%	EXPERIENCE SECTION
%----------------------------------------------------------------------------------------
% \begin{rSection}{Experience}
% %--copy and paste this region  if you need more--
% {\bf Title of the experience}{, institution associated with it} \hfill {\em start date- end date}\\
% description of the experience\\\\
% %--copy and paste this region  if you need more--
% \end{rSection}

\begin{rSection}{My Interests and Publications}
  I am now interested in \emph{theoretical computer science}, especially in \emph{cryptography}. I am fond of using mathematical methods to design or analyze an algorithm or a protocol.

{\bf \href{https://dl.acm.org/doi/10.1007/978-3-032-07891-9_8}{Two-Factor Authenticated Key Exchange with Enhanced Security from Post-Quantum Assumptions}}
\\ \emph{\textbf{Qijia Fan}, Chenhao Bao, Xuanyu Shi, Shuai Han, and Shengli Liu}. \emph{ESORICS 2025}.
\\ We proposed a two-factor authenticated key exchange protocol that provides perfect forward security and mutual explicit authentication while protecting the privacy of biometrics and passwords.

{\bf \href{https://eprint.iacr.org/2026/2482}{Edabits Generation with Quasi-Linear Communication Complexity}}
\\ \emph{Alexander Bienstock, Daniel Escudero, \textbf{Qijia Fan}, Fengrun Liu, and Yifan Song}. \emph{IACR ePrint 2026/2482}.
\\ We develop an edaBits generation protocol with quasi-linear communication complexity, facilitating efficient conversions between arithmetic and binary computation modes in secure multi-party computation.
\end{rSection}

\begin{rSection}{Awards and Scholarships}
  \begin{itemize}[leftmargin=*, parsep=0.5ex]
    \item {\bf National Scholarship}{, Ministry of Education of the People's Republic of China} \hfill{\em 2023, 2024}
    \item {\bf Outstanding Scholarship for Undergraduates, Grade A}{ (the highest level, top 5\% of the school)} \hfill{\em 2023, 2024}
    \item {\bf Irving T. Ho Memorial Scholarship} \hfill{\em 2022}
    \item {\bf Zhiyuan Hanying Ju Hua Scholarship} \hfill{\em 2023}
  \end{itemize}
\end{rSection}

%--------------------------------------------------------------------------------
%    ACTIVITIES
%-----------------------------------------------------------------------------------------------
% \newpage
\begin{rSection}{Teaching Experience}
  %--copy and paste this region  if you need more--
  {\bf Modern Algebra} \hfill {\em Spring, 2024}\\
  \emph{Teaching Assistant} with Prof. Shengli Liu
  \\ I helped students to get a better understanding of group, ring and field, and correct their homework. I also helped prepare exam questions.

  {\bf Principle and Practice of Computer Algorithms} \hfill {\em Summer, 2023}\\
  \emph{Teaching Assistant}
  \\ I helped students to understand the Raft Consensus Algorithm and implement the algorithm.
  
  {\bf Data Structures} \hfill {\em Spring, 2023}\\
  \emph{Teaching Assistant} with Prof. Huiyu Weng
  \\ I helped students to implement mergeable heaps and analyze the time complexity of some self-adjusting heaps (like skew heaps, pairing heaps, Fibonacci heaps, etc.) with amortized analysis method.
  %--copy and paste this region  if you need more--
\end{rSection}
  
%----------------------------------------------------------------------------------------
%	SKILLS SECTION
%----------------------------------------------------------------------------------------
\begin{rSection}{Skills}
  {\bf Languages: } C/C++, Go, Java, Verilog, Coq, \LaTeX, Python.
  
  \end{rSection}
  


% \begin{rSection}{Extracurricular Activities}
%   %--copy and paste this region  if you need more--
%   {\bf Shanghai Jiao Tong University Student's Choir} \hfill {\em Sept. 2023 -- Jun. 2025}\\
%   \emph{Tenor}
%   \\ I have participated in numerous campus-level performances, including the 128th Anniversary Celebration of SJTU, the opening performance of the musical \emph{Lu Xun} (the name of a famous Chinese writer), etc.
%   %--copy and paste this region  if you need more--
% \end{rSection}


\end{document}----------------------------

